\scanpage{24}

\uline{Proposition}. Soient $k$ un corps, $K$ une extension telle que $\Omega^{\unsure{1|\vphantom{1}}}_{K/k}$ soit de dim\textsuperscript{n} infinie sur $K$, et $A\inserted{=K[[t_1,\ldots,t_n]]}$ une algèbre de séries formelles sur $K$, avec $n>0$. Alors $A$, muni de sa topologie \uline{discrète}, n'est pas f\textsuperscript{t} simple sur $K$ relativement à $k$, et à fortiori n'est pas f\textsuperscript{t} simple sur $\unsure{k|K}$.

Il faut prouver que l'hom

\[ \Omega^1_{K/k}\otimes_K A\unsure{-|\longrightarrow}\Omega^1_{A/k} \]

n'a pas d'inverse à gauche, i.e. qu'il n'existe pas de dérivation

\[ D:A\unsure{-|\longrightarrow}\Omega^1_{K/k}\otimes_K A \]

qui sur $K$ induise $d_{K/k}\otimes_K\deleted{\unsure{A|1}}1_{\unsure{A|A.}}$. Or

Si on avait une telle $D$, alors les $D(t_i)$ seraient dans un sous-module $V\otimes_K A$, $V\subset\Omega_{K/k}$ de dim finie, donc il existerait une extension de \uline{type fini} $k'$ de $k$, telle que les $Dt_i$ soient dans $\operatorname{Im}\Omega_{\unsure{K'|k'}/k}\otimes_{K'}A$, donc leurs images dans $\Omega_{K/k'}\otimes_K A$ seraient \unsure{nulle|nulles}. Donc [quitte à remplacer $k$ par $k'$, ce qui ne change pas l'hyp.:] on peut supposer

% letterhead
\deleted{M. DUFOUR\\
VARITYPISTE\\
131, RUE LECOURBE\\
PARIS XV\textsuperscript{e}\\
$\blacksquare$\\
C. C. P. Paris 14.143-34\\
R. M. 140.029}
% archive sheet number: 23
